% !TEX program = xelatex
\documentclass{resume}

\usepackage[left=0.5in,top=0.45in,right=0.5in,bottom=0.25in]{geometry}
\usepackage[UTF8,fontset=none]{ctex}
\setCJKmainfont[AutoFakeBold=2.5,AutoFakeSlant=0.2]{AR PL UMing CN}
\setCJKsansfont[AutoFakeBold=2.5,AutoFakeSlant=0.2]{AR PL UMing CN}
\setCJKmonofont{AR PL UMing CN}
\usepackage{hyperref}
\usepackage{fancyhdr}
\usepackage{setspace}

\setstretch{0.95}
\renewcommand{\headrulewidth}{0pt}
\fancyhead{}
\pagestyle{fancy}
\hypersetup{colorlinks=true,urlcolor=blue}

\name{杨彦隽}
\address{\href{mailto:yanjun.yang@ed.ac.uk}{yanjun.yang@ed.ac.uk}\quad
  \href{https://blog.yanaerobe.top}{blog.yanaerobe.top}}
\address{+86 150 7499 5687\quad +44 7960 011793}

\begin{document}

\begin{rSection}{技术概述}
  专用加速器与处理器微架构博士候选人；聚焦定制 ISA、流水处理核、双可寻址存储器、
  FPGA/ASIC 原型及软硬件协同，具备从架构定义、RTL 实现到验证与性能分析的完整经验。
\end{rSection}

\begin{rSection}{教育背景}
  \begin{rSubsection}{The University of Edinburgh}{2022年9月 -- 预计2026年11月}
    {PhD Candidate in Engineering, IMNS}{英国爱丁堡}
    \item 研究方向：面向符号智能与图工作负载的专用硬件架构；工程学院全额奖学金
  \end{rSubsection}
  \begin{rSubsectionNoItems}{Nanyang Technological University}{2021年8月 -- 2022年8月}
    {MSc (Electronics)}{新加坡}
  \end{rSubsectionNoItems}
  \begin{rSubsectionNoItems}{同济大学}{2017年9月 -- 2021年7月}
    {电子科学与技术 工学学士}{中国上海}
  \end{rSubsectionNoItems}
\end{rSection}

\begin{rSection}{核心项目与研究经历}
  \begin{rSubsection}{ASOCA3：FPGA 图数据库加速系统}{2024年1月 -- 至今}
    {The University of Edinburgh}{博士核心课题}
    \item 主导系统架构与 ASOCA64 定制 ISA，定义 44-bit flit、模块化指令封装、寄存器模型及处理核到 NoC 的流量与握手接口
    \item 独立设计 supercluster 处理核、六级执行流水与控制架构，实现 valid/ready 背压、冒险检查、stall/flush 及 16 项双读单写寄存器堆
    \item 设计独立 RECONS 重构引擎，对跨流水级的复合/谓词操作进行 late binding 与路由分类，支持反馈执行、空结果抑制及跨核转发
  \end{rSubsection}

  \begin{rSubsection}{ENGINE：可重构 FPGA 双可寻址存储器}{2024年 -- 2025年}
    {The University of Edinburgh}{独立设计与实现}
    \item 提出并完成 BRAM-based dually-addressable memory 的架构、RTL、FPGA 实现与评估，以 array/page/RAMB 映射同时支持地址与内容寻址且无需复制数据
    \item 在 Virtex-7 上以 4 BRAM36、683 LUT、248 FF 实现 $8\times512\times36$ 存储，达到 129.27 MHz 与 100\% BRAM 存储效率；在 UltraScale+ 上验证 302 实例、约 124 万 entries
  \end{rSubsection}

  \newpage

  \begin{rSubsection}{ASOCA2：双可寻址近存图加速 ASIC}{2022年11月 -- 2023年12月}
    {The University of Edinburgh}{博士子课题}
    \item 主导 \textit{Views} 数据映射、ISA 与 hypercluster 架构，参与 CTRL、BUF、ONECON、LOI、COMP 等数字外设的多版 RTL、集成与验证
    \item 参与 180 nm 原型芯片综合与流片交付，主导顶层验证及硅后测试与故障分析；硅前 48 项指令/初始化检查全部通过，硅后发现数字 I/O/scan 边界信号衰减
  \end{rSubsection}

  \begin{rSubsection}{AIDE：架构开发、验证与评估基础设施}{2025年11月 -- 至今}
    {The University of Edinburgh}{博士子课题}
    \item 独立开发 Makefile 自动化设计流、ASOCA64 assembler/simulator，并扩展 Python 行为模型以覆盖 traffic、register 与 instruction reconstruction
    \item 主导定制 compliance 环境，以 sequencer、monitor、scoreboard 和 functional coverage 验证 RTL；设计 Neo4j 对比方法与工作负载分析流程
  \end{rSubsection}

  \begin{rSubsection}{CoNM：四级流水 RV32I 处理器}{2021年3月 -- 2021年6月}
    {同济大学}{本科毕业设计}
    \item 使用 Verilog 设计带静态分支预测的四级流水 RV32I CPU 内核
    \item 在 PYNQ-Z1 FPGA 开发板完成综合实现与功能验证
  \end{rSubsection}
\end{rSection}

\begin{rSection}{专业技能}
  \begin{rSubsection}{体系结构}{}{}{}
    \item Processor/accelerator microarchitecture、ISA、pipeline、memory architecture、NoC traffic/interface、hardware-software co-design
  \end{rSubsection}
  \begin{rSubsection}{硬件与验证}{}{}{}
    \item SystemVerilog/Verilog、RTL design、assertion-based verification、custom compliance、FPGA/ASIC prototyping
  \end{rSubsection}
  \begin{rSubsection}{工具与编程}{}{}{}
    \item Vivado、VCS、ModelSim；Python、C/C++、Bash、Tcl、Git
  \end{rSubsection}
\end{rSection}

\begin{rSection}{代表性论文}
  \begin{rSubsectionNoItems}{\textit{Views}: A Hardware-friendly Graph Database Model For Storing Semantic Information}
    {2025}{arXiv preprint}{Graph representation / accelerator architecture}
  \end{rSubsectionNoItems}
  \begin{rSubsectionNoItems}{A Resource-efficient Dually-addressable Memory Architecture on FPGA}
    {2025}{ISCAS}{Dually-addressable memory / FPGA}
  \end{rSubsectionNoItems}
  \begin{rSubsectionNoItems}{A Modular Graph Database Accelerator}
    {2024}{ECML PKDD}{Associative-memory accelerator / ASIC}
  \end{rSubsectionNoItems}
\end{rSection}

\end{document}
